% !TeX root = ../samplepaper.tex
\section{Introduction}
Focused Ant Colony Optimization（FACO）通过复用参考路径的大部分结构，减少旅行商问题（TSP）中重复构造完整路径的开销\cite{skinderowicz2022}。每只蚂蚁只主动引入有限数量的新连接，再由local search改善所得路径。该机制将搜索集中在已有优质解附近，也带来一个控制问题：搜索应从哪里开始、修改多少结构，以及何时转向另一条参考路径？这些决策共同决定了搜索范围和计算工作量。

本文研究利用Genetic Programming（GP）学习上述控制决策。其动机是，固定配置难以区分不同搜索状态：小幅修改可能反复返回参考路径，较大修改则可能增加无效搜索；更换参考路径提供了另一种探索方式，但也会改变可供利用的结构。GP能够以短表达式组合状态与候选动作的特征，形成可在求解过程中反复调用的priority function。离线进化根据完整FACO求解的质量评价这些规则，在线求解则由规则选择动作，蚁群继续负责路径构造与local search。

围绕这一问题，本文定义参考路径切换、起点区域和扰动强度的联合控制接口，并研究single-tree联合评分与multi-tree条件决策两种实现。实验在TSP500上学习策略，通过原生FACO及共享GPU底座的静态策略评价收益，再考察TSP1000迁移、进化规则和计算成本。强度较大的静态策略作为必要对照，用于判断学习所得规则是否超过了简单的参数选择。

\section{Background and Related Work}
\subsection{FACO中的结构复用与控制}
蚁群优化通过信息素积累搜索经验，MAX--MIN Ant System以信息素上下界调节探索与利用\cite{stuetzle2000}。FACO进一步利用参考路径，只主动改变部分连接\cite{skinderowicz2022}。令$q$为参考路径，$P_a$为动作$a$控制的构造过程，$L$为local search，一次候选解生成可表示为
\begin{equation}
 q\xrightarrow{P_a}y\xrightarrow{L}z.
\end{equation}
控制的作用体现在$P_a$生成什么样的候选路径，以及这些候选在$L$之后是否有助于继续改善解。参考路径切换改变$q$，区域选择影响修改起点，MNE（minimum number of new edges）控制主动引入新连接的目标数。MNE不等于最终路径与$q$的边差异数。

后续自适应FACO研究了在线参数调节\cite{skinderowicz2024}。本文的GPU底座采用该后续实现的节点重定位式构造：复制参考路径后，依据候选邻域中的信息素与距离启发式选择节点，再将节点移动到当前节点之后。GP在此基础上学习三项控制的组合。2022原生FACO与该底座的构造和初始化不同，因此实验同时保留原生方法及共享底座的静态对照。

\subsection{GP学习决策规则}
GP hyper-heuristic以启发式规则为搜索对象\cite{burke2013}。一个individual表示priority function；求解过程使用该函数为合法候选评分，并执行排名最高的选择。其fitness来自应用该规则后生成的解。例如，orienteering中的GP对候选访问点评分\cite{mei2018orienteering}，UCARP中的GP学习车辆的任务选择规则\cite{maclachlan2020}。这类方法将规则学习与解的构造连接起来，也便于从表达式分析候选偏好。

当一个求解过程包含多个决策时，multi-tree representation可将不同角色的规则放入同一个individual。Zhang等联合进化routing与sequencing规则，并提出tree-swapping crossover\cite{zhang2018multitree}；Sun等以三棵树分别选择task、cloud和resource，通过共同fitness评价三者的配合\cite{sun2024multicloud}。本文比较完整动作的联合评分与三个控制项的条件选择。二者表达相同类型的控制任务，但具有不同的搜索空间、特征可见性和繁殖方式，实验评价的是这两套完整设计。
